\section{방법}

\subsection{온톨로지 기저}

모든 실험은 tool catalog로부터 구성한 per-head ontology basis $B_{\mathrm{ont}} \in \mathbb{R}^{L \times H_{\mathrm{kv}} \times d \times r}$를 사용한다. 직관적으로 이 기저는 금융, 뉴스, 여행, 미디어 같은 tool facet에서 반복적으로 나타나는 방향을 span한다. 평가 시점에서 기저는 고정되며, query-side와 key-side 개입이 동일한 기저를 공유한다.

\subsection{실용적인 query-space coverage bias}

query tensor $q$에 대해 구현된 query-side 개입은 다음과 같다.
\[
q' = q + \beta B_{\mathrm{ont}} B_{\mathrm{ont}}^\top q.
\]
$\beta < 0$이면 연산자는 ontology subspace에 정렬된 query 성분을 빼준다. 코드에서는 이를 \texttt{q\_proj}에 대한 forward hook으로 구현한다. 이 방법은 history-free다. 어떤 도구가 이미 출력되었는지를 명시적으로 추적하지 않는다. 따라서 우리는 이를 더 강한 이상적 coverage operator의 실용적 근사로 해석하며, 문자 그대로의 per-step facet-state tracking으로 해석하지 않는다.

\subsection{정적인 key-side bias}

비교 대상으로는 동일한 기저를 사용하는 key-side 개입을 둔다.
\[
k' = k + \alpha B_{\mathrm{ont}} B_{\mathrm{ont}}^\top k.
\]
이 개입은 정의상 stationary하다. 모든 디코딩 단계에서 같은 ontology-aligned key direction을 증폭한다. 첫 번째로 출력된 도구가 이미 지배적 facet을 차지하고 있다면, 이 연산자는 이후 단계에서 그 facet을 억제할 방법이 없다.

\subsection{작업 가설}

본 논문의 작업 가설은 두 부분으로 이루어진다. 첫 번째는 \emph{정확도 가설}이다. 음의 query-side ontology bias는 지배적인 facet의 반복 사용을 줄여주므로 multi-tool recovery를 향상시켜야 한다. 두 번째는 \emph{특이성 가설}이다. ontology basis를 동일 rank의 random basis나 feature-shuffled basis로 대체하면 효과가 사라져야 한다. 실험은 이 두 주장을 직접 검증하도록 설계되어 있다.

